\section*{Chapitre \ref{ch:ml:monitoring-and-retraining} --- Surveillance et réentraînement}

\begin{solution}{exo:ml:monitoring-and-retraining:1}
Entrées~: la dispersion croissante de l'érosion lente (vue par l'\omterm{def:ml:monitoring-and-retraining:psi}{indice de stabilité de la population} (PSI) des variables et par la statistique de Kolmogorov--Smirnov, et par rien d'autre). Sorties~:
le changement d'unité, qui se lit comme une hausse de la part située dans le seuil de négociation, de 50~\% à 63~\%, que les entrées voient aussi, mais que seules
les sorties traduisent en négociation (un quart de valeurs négociées en moins). Résultats~: l'inversion, invisible pour les entrées et les sorties, parce que
ni les variables ni la distribution des prédictions ne changent.
\end{solution}

\begin{solution}{exo:ml:monitoring-and-retraining:2}
Avec 200 valeurs, chaque décile en compte 20, et le seul bruit d'échantillonnage donne à l'indice quotidien de chaque variable une dispersion telle que le
plus grand des cinq dépasse 0.10 plus d'un jour sur vingt. La règle empirique a été fixée pour de grands portefeuilles de crédit et ne dit rien
des taux d'erreur~; un seuil de 0.10 ici déclencherait une alerte plus d'une fois par mois, et l'équipe apprendrait à l'ignorer.
\end{solution}

\begin{solution}{exo:ml:monitoring-and-retraining:3}
Ses fausses alertes arrivent au hasard, environ un jour sur 21, si bien que l'attente de la première à partir de n'importe quelle date suit une loi géométrique de médiane
d'environ 14 jours (la colonne «~aucune~» du tableau indique 10 à 17 jours pour les cinq moniteurs). Une première alarme n'a de sens que lorsqu'elle
intervient bien avant cela.
\end{solution}

\begin{solution}{exo:ml:monitoring-and-retraining:4}
$\Delta=0.103-0.087=0.017$ (valeur non arrondie 0.0168), $\sigma=0.071$, $z_{0.95}+z_{0.90}=1.645+1.282=2.927$~: $n=(2.927\times0.071/0.0168)^2
\approx150$ jours. L'écart-type du coefficient quotidien décroît à peu près comme l'inverse de la racine carrée du nombre de valeurs~;
avec 2\,000 valeurs il vaut environ $0.071/\sqrt{10}=0.022$, et $n$ est divisé par dix, soit environ 15 jours.
\end{solution}

\begin{solution}{exo:ml:monitoring-and-retraining:5}
L'erreur type du test est estimée à partir des jours observés jusqu'ici. Sur trois jours, l'écart-type empirique des différences
peut être minuscule par hasard, ainsi que l'erreur type qui en découle, et une différence moyenne modeste paraît alors écrasante~: avec les verdicts
issus du troisième jour, l'une des simulations du chapitre affiche une p-valeur de $1.3\times10^{-6}$ ce jour-là. La garantie d'erreur du test
suppose une erreur type connue~; dix jours constituent une période de chauffe (burn-in) qui rend l'estimation utilisable. Sans défaillance, les verdicts issus du troisième
jour promeuvent 16 challengers sans effet sur 100, contre 2 avec la période de chauffe.
\end{solution}

\begin{solution}{exo:ml:monitoring-and-retraining:6}
L'alarme était juste, mais le correctif était faux. Le changement du fournisseur était une erreur de données~; le réentraînement a appris au modèle à multiplier la
variable en lots par cent fois le coefficient, enfouissant l'erreur dans le modèle. Quand le fournisseur corrige le flux de données, le
modèle réentraîné est faux d'un facteur cent sur cette variable, et l'alarme se déclenche de nouveau. Il faut corriger les données (convertir les unités à la
frontière du système, ajouter un contrôle de schéma ou de plage) et conserver le modèle.
\end{solution}

\begin{solution}{exo:ml:monitoring-and-retraining:7}
\texttt{ml\_monitor.pooled}~: le seuil de l'indice agrégé pour une alerte par mois est de 0.0056, vingt fois plus bas que le seuil quotidien, car
4\,000 échantillons par variable portent beaucoup moins de bruit que 200. Des fenêtres consécutives partagent 19 de leurs 20 jours, si bien que la
statistique évolue lentement, et ses fausses alarmes arrivent par épisodes, soit 3.7 jours d'alarme par alerte en moyenne. La première alarme ne fait pas mieux qu'un moniteur aveugle
(médiane de 15.5 jours pour l'érosion, 18 sans défaillance), et la règle cinq sur 21 est déclenchée par un seul faux épisode (médiane de 20
jours sans défaillance), si bien qu'aucun de ces deux décomptes ne signifie plus ce qu'il signifiait. Mesuré par un mois d'alarme continue, l'indice agrégé voit
l'érosion après une médiane de 22.5 jours, contre 120 pour l'indice quotidien~; mais quatre des 20 simulations sans défaillance montrent aussi un faux
mois d'alarme continue. L'agrégation gagne en sensibilité aux dérives lentes, au prix d'alarmes moins nombreuses mais plus longues, et plus difficiles à écarter.
\end{solution}

\begin{solution}{exo:ml:monitoring-and-retraining:8}
Aucun moniteur d'entrées ou de sorties ne peut le voir~: la distribution des variables est inchangée, et il en va de même, au premier ordre, pour celle des
prédictions. Seuls les résultats le peuvent, et la perte est progressive~: le coefficient tombe à 0.081 à la fin de l'érosion, et la formule du chapitre
demande environ 80 jours à ce niveau, si bien que l'alarme n'arriverait que des mois après le début de l'érosion, si tant est qu'elle arrive, à raison d'une fausse alerte par
mois. La réponse pratique n'est pas un moniteur mais un calendrier~: réestimer ou revalider régulièrement le modèle (chapitre~12), comparer chaque
variable~: sa contribution récente (son coefficient réestimé sur une fenêtre glissante) avec celle de l'entraînement, et examiner l'attribution de
l'IC réalisé par variable, ce qui concentre les éléments probants sur la variable qui s'érode.
\end{solution}

\begin{solution}{pb:ml:monitoring-and-retraining:1}
\textbf{Partie I.}
\begin{enumerate}
\item Deux cents valeurs par jour, cinq variables standardisées, des rendements du lendemain avec des coefficients de 0.08, $-0.06$ et 0.05 sur les trois premières
plus un bruit unitaire~; un modèle linéaire estimé sur 60 jours tourne sur 400 jours, en négociant la moitié des valeurs aux prévisions les plus élevées. À partir du jour
200~: la deuxième variable cent fois plus petite~; l'effet de la première variable s'inverse~; celui de la troisième s'estompe jusqu'à zéro sur 150
jours, tandis que sa dispersion croît de moitié.
\item Entrées (distributions des variables face à la référence d'entraînement~; immédiatement), sorties (distribution des prédictions, part située dans
le seuil de négociation~; immédiatement), résultats (coefficient d'information réalisé, P\&L~; après l'horizon de prévision, et bruités).
\item $\sum_b(q_b-r_b)\ln(q_b/r_b)$ sur les déciles de la référence~; seuils 0.115 (plus grand indice des variables), 0.115 (la plus grande
statistique de Kolmogorov--Smirnov), 0.089 (indice des prédictions), 0.07 (variation de la part située dans le seuil).
\item Parce qu'une personne n'est alertée qu'une fois par épisode, et non une fois par jour de celui-ci~; calibrer sur les jours d'alarme ferait
paraître les statistiques persistantes plus bruitées qu'elles ne le sont.
\end{enumerate}
\textbf{Partie II.}
\begin{enumerate}[resume]
\item Unilatéral, sur le déficit du coefficient quotidien par rapport à sa valeur attendue de 0.103, tolérance $k=0.035$ (la moitié de l'écart-type
quotidien de 0.071), seuil $h=0.11$ choisi sur des simulations sans défaillance pour une alerte par mois, réinitialisé après chaque alerte.
\item Environ $((z_{0.95}+z_{0.90})\sigma/\Delta)^2$~: 4 jours pour l'inversion, environ 150 pour le changement d'unité, environ 80 pour l'érosion à son
terme.
\item Le changement d'unité retire la contribution d'une variable à la prévision~; le coefficient ne tombe que de 0.103 à 0.087, une
baisse égale au quart du bruit d'une journée.
\item Une première alarme est la première alerte après la défaillance, que produit aussi un moniteur aveugle dans les deux semaines environ~; une alarme soutenue
est cinq jours d'alarme sur 21, seuil qu'aucun moniteur n'atteint sans défaillance.
\end{enumerate}
\textbf{Partie III.}
\begin{enumerate}[resume]
\item Les quatre moniteurs d'entrées et de sorties alertent dès le premier jour, soutenus d'emblée~; le moniteur de résultats alerte une fois après une médiane de
9.5 jours et ne persiste jamais.
\item Seul le CUSUM sur l'IC~: en deux jours dans 13 des 20 simulations, soutenu d'emblée.
\item Rien ne voit la précision perdue~; les moniteurs d'entrées voient la dispersion croissante, soutenue après une médiane de 42 jours (indice des
variables) et 60 (Kolmogorov--Smirnov).
\item Rien dans le tableau~; une revalidation programmée, ou un moniteur de la contribution de chaque variable au coefficient réalisé.
\end{enumerate}
\textbf{Partie IV.}
\begin{enumerate}[resume]
\item \emph{Le mardi où le modèle s'est tu.} Avec une fausse alerte par mois (médianes sur 20 simulations~; première alarme / soutenue)~:
changement d'unité, chaque moniteur d'entrées et de sorties 0 / 0 jour, CUSUM sur l'IC 9.5 jours / jamais~; inversion, moniteurs d'entrées et de sorties pas meilleurs
qu'un moniteur aveugle (10 à 17 jours, jamais soutenus), CUSUM sur l'IC 1 / 0 jour~; érosion lente, indice des variables 12 / 42 jours, Kolmogorov--Smirnov 8 /
60, indice des prédictions, part dans le seuil et CUSUM sur l'IC, jamais soutenus.
\item Un challenger réestimé sur 40 jours postérieurs à l'inversion (coefficient d'information 0.105 contre 0.003 pour le champion) est promu dans
les 20 simulations après une médiane de 10 jours de \omterm{def:ml:monitoring-and-retraining:shadow}{déploiement fantôme}, au plus 19~; sans défaillance, 2 challengers sur 100 sont promus.
\item Un challenger affecté d'une erreur de signe est annulé par le CUSUM après une médiane de 3 jours~; le coefficient d'information abandonné
s'élève à 0.26 jour-IC à pleine taille et à 0.026 au dixième.
\item Alerter le propriétaire des données sur les alarmes de part dans le seuil et de variables, arrêter ou réduire le modèle, trouver le changement d'unité et corriger le flux de données~;
et non réentraîner.
\item Réentraîner après un changement confirmé de la relation, dès qu'il y a assez de données postérieures au changement, et selon un calendrier~; retirer le modèle lorsque
les challengers réestimés ne battent plus une référence nulle.
\item Alarmes d'entrées~: le propriétaire des données~; alarmes de sorties~: le desk et le propriétaire du modèle~; alarmes de résultats~: le propriétaire du modèle et le risque.
\item Chaque alarme, la personne qui en a accusé réception, l'action (correction des données, réduction de taille, retour arrière, promotion) avec sa raison, rattachée au
modèle dans le registre.
\item Un taux de fausses alarmes auquel l'équipe saura répondre et un délai mesuré sur les défaillances qu'il est censé détecter.
\end{enumerate}
\end{solution}

\begin{solution}{iq:ml:monitoring-and-retraining:1}
Entrées (schéma, fraîcheur, distributions des variables), sorties (distribution des prédictions, positions, taux de rotation, part des transactions) et
résultats (IC réalisé, P\&L, mark-outs). Chaque couche détecte des défaillances que les autres manquent~: un changement d'unité se voit fortement dans les entrées et faiblement dans
les résultats~; un changement de régime est muet dans les entrées et flagrant dans les résultats.
\iqlookfor{trois couches, avec un exemple pour chacune.}
\end{solution}

\begin{solution}{iq:ml:monitoring-and-retraining:2}
$\sum_b(q_b-r_b)\ln(q_b/r_b)$ sur les classes de la référence. Choisir le seuil pour un taux de fausses alarmes sur des données sans défaillances
(historique ou simulation), avec la taille d'échantillon réelle et le nombre de variables surveillées~; les règles empiriques de 0.10 et 0.25 ignorent
ces deux éléments.
\iqlookfor{la formule et la calibration à un taux d'erreur.}
\end{solution}

\begin{solution}{iq:ml:monitoring-and-retraining:3}
$\Delta=0.02$, $\sigma=0.07$~: $((1.645+1.282)\times0.07/0.02)^2\approx105$ jours, soit cinq mois, pour un test unilatéral à 5~\% avec 90~\%
de puissance~; davantage de titres le raccourcit, car $\sigma$ décroît à peu près comme la racine carrée de leur nombre.
\iqlookfor{le calcul de la taille d'échantillon et sa dépendance à la largeur (breadth).}
\end{solution}

\begin{solution}{iq:ml:monitoring-and-retraining:4}
\omterm{def:ml:monitoring-and-retraining:shadow}{Déploiement fantôme}~: le nouveau modèle reçoit les entrées en direct et ses décisions sont évaluées mais non exécutées~; adapté à la comparaison des prévisions, sans risque.
Déploiement canari~: il négocie une petite part du capital~; adapté lorsque ses transactions modifient les résultats (impact, exécutions). \omterm{def:ml:monitoring-and-retraining:shadow}{Champion--challenger}~: la
discipline qui encadre les deux, garder le modèle en place jusqu'à ce que le challenger remporte un test choisi à l'avance.
\iqlookfor{ce que mesure chacun et les risques qu'il comporte.}
\end{solution}

\begin{solution}{iq:ml:monitoring-and-retraining:5}
Pas en premier. Vérifier si la dérive relève d'une erreur de données (unités, schéma, changement de fournisseur) ou d'un véritable changement dans la population~; corriger les
erreurs à la source. Si la population a réellement changé, vérifier les résultats~: un modèle peut rester robuste à un \omterm{def:ml:online-learning-and-drift:drift}{décalage des covariables} (covariate shift)~; réentraîner
lorsque les résultats ou une nouvelle validation montrent qu'il ne l'est pas, et promouvoir par un \omterm{def:ml:monitoring-and-retraining:shadow}{déploiement fantôme}.
\iqlookfor{d'abord le chemin des données, puis les résultats avant le réentraînement.}
\end{solution}

\begin{solution}{iq:ml:monitoring-and-retraining:6}
Mesurer le taux de fausses alarmes de chaque dispositif de surveillance et son délai sur les défaillances qu'il est censé détecter~; retirer les dispositifs qui ne font jamais mieux qu'un contrôle aveugle~;
recalibrer les autres à un taux auquel l'équipe saura y répondre (en comptant les alertes à leur déclenchement)~; acheminer chaque alarme vers son propriétaire~; regrouper les alarmes liées~;
et couvrir les pertes lentes par des revues programmées, plutôt que par des alertes.
\iqlookfor{calibration à un taux, valeur mesurée pour chaque dispositif, acheminement.}
\end{solution}
